\documentclass[10pt]{article}
\usepackage[T1]{fontenc}
\usepackage[utf8]{inputenc}
\usepackage[french]{babel}
\usepackage{amsmath,amssymb,mathrsfs,mathpazo}
\usepackage{geometry}
\geometry{a4paper,margin=2cm}
\usepackage{graphicx,array,tabularx,tabularray,colortbl,multirow,multicol,booktabs,enumitem,pifont,ulem,fancybox,diagbox,cancel}
\usepackage{pstricks,pstricks-add,pst-plot,pst-tree,pst-node,pst-text,pst-func,pst-math,pst-bspline,pst-3dplot,pst-eucl,pst-solides3d}
\let\pscalculate\undefined
\usepackage{pst-bezier}
\usepackage{tikz,pgfplots,tkz-tab}
\usetikzlibrary{arrows,arrows.meta,patterns,calc,positioning,trees,decorations.markings,shapes.geometric,shapes.misc,shapes.arrows,intersections,decorations.pathmorphing,decorations.pathreplacing,backgrounds,shadows,babel}
\pgfplotsset{compat=1.18}
\usepgfplotslibrary{fillbetween}
\usepackage[np]{numprint}
\usepackage{listings,makecell,pas-tableur,twoopt}
\let\barre\undefined
\usepackage[pstricks]{tablvar}
\definecolor{darkgreen}{RGB}{0,100,0}
\providecommand{\red}{\color{red}}
\providecommand{\blue}{\color{blue}}
\providecommand{\green}{\color{green}}
\providecommand{\black}{\color{black}}
\providecommand{\white}{\color{white}}
\providecommand{\gray}{\color{gray}}
\providecommand{\cyan}{\color{cyan}}
\providecommand{\magenta}{\color{magenta}}
\providecommand{\yellow}{\color{yellow}}
\providecommand{\violet}{\color{violet}}
\providecommand{\orange}{\color{orange}}
\providecommand{\euro}{\texteuro}
\providecommand{\EUR}{\texteuro}
\providecommand{\ieme}{\textsuperscript{e}}
\providecommand{\degres}{\ensuremath{^\circ}}
\providecommand{\degree}{\ensuremath{^\circ}}
\providecommand{\vv}[1]{\overrightarrow{#1}}
\providecommand{\Coord}[1]{#1}
\providecommand{\decosix}{}
\providecommand{\starredbullet}{\ensuremath{\bullet}}
\pagestyle{empty}
\setlength{\parindent}{0pt}
\usepackage{adjustbox}\DeclareUnicodeCharacter{2236}{:}\DeclareUnicodeCharacter{2190}{\ensuremath{\leftarrow}}
\providecommand{\R}{\mathbb{R}}
\providecommand{\N}{\mathbb{N}}
\providecommand{\Z}{\mathbb{Z}}
\providecommand{\Q}{\mathbb{Q}}
\providecommand{\D}{\mathbb{D}}
\providecommand{\C}{\mathbb{C}}
\providecommand{\e}{{\mathrm{e}}}
\providecommand{\vect}[1]{\overrightarrow{#1}}
\providecommand{\vectt}[1]{\overrightarrow{\mathrm{#1}}}
\providecommand{\vvt}[1]{\overrightarrow{\mathrm{#1}}}
\providecommand{\np}[2][]{{#2~#1}}
\providecommand{\vv}[1]{{\overrightarrow{#1}}}
\renewcommand{\d}{{\mathrm{d}}}
\providecommand{\degres}{{^\circ}}
\providecommand{\degre}{{^\circ}}
\providecommand{\blue}{\color{blue}}
\providecommand{\red}{\color{red}}
\providecommand{\enskip}{\quad}
\providecommand{\llbracket}{\lbrack\!\lbrack}
\providecommand{\rrbracket}{\rbrack\!\rbrack}
\providecommand{\euro}{{\text{€}}}
\providecommand{\no}{{\text{n°}}}
\providecommand{\interval}[2]{{[#1;#2]}}
\providecommand{\intervFO}[2]{{[#1;#2[}}
\providecommand{\dotfill}{\ldots}

\providecommand{\ds}{\displaystyle}
\providecommand{\textding}[1]{#1}
\providecommand{\Oij}{\(\mathrm{O};\vec{\imath},\vec{\jmath}\)}
\providecommand{\Oijk}{\(\mathrm{O};\vec{\imath},\vec{\jmath},\vec{k}\)}
\providecommand{\Ouv}{\(\mathrm{O};\vec{u},\vec{v}\)}
\providecommand{\intervFF}[2]{\ensuremath{[#1;#2]}}
\renewcommand{\euro}{\ifmmode\text{\texteuro}\else\texteuro\fi}
\DeclareUnicodeCharacter{00B0}{\ifmmode^{\circ}\else\textdegree\fi}
\DeclareUnicodeCharacter{20AC}{\ifmmode\text{\texteuro}\else\texteuro\fi}
% Énoncés : Éducation nationale. Transcription et corrigés : APMEP, Denis Vergès et contributeurs. Découpage et classement : CoopMaths, Sébastien Lozano, Éric Elter et contributeurs. Adaptation de la présentation : Maths Explorer.
% CC BY-SA : https://coopmaths.fr/www/about/
% https://coopmaths.fr/alea/static/bac/2025/tex/bac_2025_06_sujet2_etrangers_3.tex
\begin{document}
\noindent\textit{APMEP / CoopMaths. CC BY-SA. Présentation adaptée par Maths Explorer.}\par\medskip
Le codage \og base64 \fg, utilisé en informatique, permet de représenter et de transmettre des messages et d'autres données telles que des images, en utilisant 64 caractères: les 26 lettres majuscules, les 26 lettres minuscules, les chiffres de 0 à 9 et deux autres caractères spéciaux.

\emph{Les parties {\rm A, B} et {\rm C} sont indépendantes.}

\medskip

\textbf{Partie A}

\medskip

Dans cette partie, on s'intéresse aux séquences de 4 caractères en base64.
Par exemple, \og gP3g \fg est une telle séquence.
Dans une séquence, l'ordre est à prendre en compte: les séquences \og m5C2 \fg et \og 5C2m \fg ne sont pas identiques.

\medskip

\begin{enumerate}
\item Déterminer le nombre de séquences possibles.
\item Déterminer le nombre de séquences si l'on impose que les 4 caractères sont
différents deux à deux.
\item 
	\begin{enumerate}
		\item Déterminer le nombre de séquences ne comportant pas de lettre A majuscule
		\item En déduire le nombre de séquences comportant au moins une lettre A majuscule.
		\item Déterminer le nombre de séquences comportant exactement une fois la lettre A majuscule.
		\item Déterminer le nombre de séquences comportant exactement deux fois la lettre A majuscule.
	\end{enumerate}
\end{enumerate}

\medskip

\textbf{Partie B}

\medskip

On s'intéresse à la transmission d'une séquence de 250 caractères d'un ordinateur à un autre. On suppose que la probabilité qu'un caractère soit mal transmis est égale à 0,01 et que les transmissions des différents caractères sont indépendantes entre elles.
On note X la variable aléatoire égale au nombre de caractères mal transmis.

\medskip

\begin{enumerate}
\item On admet que la variable aléatoire X suit la loi binomiale. Donner ses paramètres.
\item Déterminer la probabilité que tous les caractères soient bien transmis. \emph{On donnera
l'expression exacte, puis une valeur approchée à $10^{-3}$ près}.
\item Que pensez-vous de l'affirmation suivante: \og La probabilité que plus de 16
caractères soient mal transmis est négligeable\fg ?
\end{enumerate}

\medskip

\textbf{Partie C}

\medskip

On s'intéresse maintenant à la transmission de 4 séquences de 250 caractères.

On note $X_1,\: X_2,\: X_3$\: et $X_4$ les variables aléatoires correspondant aux nombres de caractères mal transmis lors de la transmission de chacune des 4 séquences. 

On admet que les variables aléatoires $X_1,\: X_2,\: X_3$\: et $X_4$ sont indépendantes entre elles et suivent la même loi que la variable aléatoire $X$ définie en partie B.

On note $S = X_1 + X_2 + X_3 + X_4.$

Déterminer, en justifiant, l'espérance et la variance de la variable aléatoire $S$.

\end{document}
